\documentclass[11pt,letterpaper]{article}
\usepackage[margin=1in]{geometry}
\usepackage{amsmath,amssymb,amsthm}
\usepackage{mathtools}
\usepackage{hyperref}
\usepackage{bm}
\usepackage{booktabs}
\usepackage{enumitem}
\usepackage{listings}
\usepackage{microtype}
\usepackage{graphicx}
\usepackage{float}

\title{\textbf{The Universal Multiplicity Constant \(\Lambda_m\)}\newline\large
Prime-Indexed Operator Dynamics, Contraction Bounds, and Governance-Integrated Validation\\
Comprehensive Report and Prior Art Defensive Publication}

\author{Citizen Gardens \\ The Foundation of Multiplicity Foundation}
\date{26 April 2026}

\begin{document}
\maketitle

\begin{abstract}
This document records a complete development cycle around the Universal
Multiplicity Constant \(\Lambda_m\), a prime-indexed operator framework,
and a Phase Mirror--style governance envelope.
We formalise a discrete evolution law
\(X_{t+1} = \Xi X_t + m T(X_t)\) with \(\Xi\) built from prime-indexed
operators, derive a sufficient contraction criterion \(F + c < 1\), and
present negative control results falsifying a ``nice interference''
hypothesis. We then outline an optimisation programme that constructs
prime-structured families achieving theorem-level contraction.

The goal is twofold: (i) provide a mathematically coherent backbone for
\(\Lambda_m\)-guided dynamics and (ii) create a timestamped prior-art
record of the architecture, methods, and governance integration.
\end{abstract}
\newpage
\tableofcontents
\newpage
\section{Executive Summary}

This report summarises a research and governance arc in which the
Universal Multiplicity Constant \(\Lambda_m\) and prime-indexed operators
are used to structure a nonlinear discrete dynamical system and its
validation protocol.

The central objects are:

\begin{itemize}[leftmargin=*]
  \item A discrete evolution law of the form
    \[
      X_{t+1} = \Xi X_t + m T(X_t),
    \]
    where \(\Xi\) is a real-valued linear operator built by aggregating
    complex prime-indexed operators, \(m\) is a small scalar coupling,
    and \(T\) is a saturating nonlinearity (e.g.\ tanh).
  \item A contraction criterion
    \[
      F + c < 1, \qquad F = \|\Xi\|_{\rm op},\quad c = |m|L_T,
    \]
    taken as the sole theorem-level sufficient condition for uniform
    contraction of the map.
  \item Negative control results in which diagonal, Fourier-conjugated,
    and permutation-conjugated prime-phase operators all produce operator
    norms \(F\) equal to the naive weight sum \(\sum |w_p|\), thereby
    falsifying an intuitive ``nice interference'' hypothesis.
  \item An explicit optimisation programme that treats the operator norm
    \(F\) itself as an objective over phases and signed/complex
    weights in prime-indexed families, producing explicit regimes in which
    \(F + c < 1\) and the contraction theorem applies.
  \item An accompanying governance envelope, expressed through
    Architecture Decision Records (ADRs), that treats falsification and
    optimisation as first-class events in a Phase Mirror--style
    governance stack.
\end{itemize}

The key claim is not that this is the most general or powerful
prime-indexed framework, but that the combination

\begin{enumerate}[leftmargin=*]
  \item prime-indexed operator families,
  \item an explicit operator norm contraction criterion,
  \item a falsifiable protocol with frozen parameters, and
  \item a governance surface that constrains what may be claimed,
\end{enumerate}

constitutes a distinctive pattern worth documenting as prior art.

\section{Prime-Indexed Operator Framework}

\subsection{State Space and Norms}

Let \(d\in\mathbb{N}\) be the dimension of the state space.
We work with vectors in \(\mathbb{R}^d\) (or \(\mathbb{C}^d\)
when convenient) equipped with the Euclidean norm
\(\|x\|_2 = (\sum_{k=1}^d |x_k|^2)^{1/2}\).
For matrices \(A\in\mathbb{C}^{d\times d}\) we use

\begin{itemize}[leftmargin=*]
  \item the operator (spectral) norm
  \[
    \|A\|_{\rm op} = \sup_{\|x\|_2=1} \|Ax\|_2 = \sigma_{\max}(A),
  \]
    where \(\sigma_{\max}(A)\) is the largest singular value; and
  \item the entrywise real part \(\operatorname{Re}(A)\) defined by
    \((\operatorname{Re}(A))_{ij} = \operatorname{Re}(A_{ij})\).
\end{itemize}

\subsection{Prime-Indexed Operator Families}

Let \(\mathcal{P}\) be a finite set of primes. For each prime \(p\in\mathcal{P}\)
we define:

\paragraph{Diagonal prime-phase operators.}
A diagonal unitary \(D_p\in\mathbb{C}^{d\times d}\) with entries
\[
  (D_p)_{kk} = e^{2\pi i k/p}, \qquad k = 0, \dots, d-1.
\]

\paragraph{Fourier-conjugated operators.}
Let \(F\in\mathbb{C}^{d\times d}\) be the unitary DFT matrix
(obeying \(F^\dagger F = I_d\)). Define
\[
  U_p^{\rm (F)} = F^\dagger D_p F.
\]

\paragraph{Permutation-conjugated operators.}
For a permutation matrix \(P_p\in\mathbb{R}^{d\times d}\) we define
\[
  U_p^{\rm (perm)} = P_p^\dagger D_p P_p.
\]

Each family \(\{U_p\}_{p\in\mathcal{P}}\) is norm-bounded with
\(\|U_p\|_{\rm op} = 1\) whenever \(D_p\) is unitary.

\subsection{Weights and Aggregation}

We introduce a weight pattern \(\{w_p\}_{p\in\mathcal{P}}\) of the form
\[
  w_p = \frac{\Omega(p)}{p^\alpha},
\]
where \(\Omega(p)\) is a multiplicity function and \(\alpha > 1\) a fixed
exponent; for primes, \(\Omega(p)=1\) in the simplest case.

Given a family \(\{U_p\}_{p\in\mathcal{P}}\) and weights \(\{w_p\}\),
we define the aggregated operator
\begin{equation}
  \Xi = \operatorname{Re}\Bigl(\sum_{p\in\mathcal{P}} w_p U_p\Bigr)
  \in \mathbb{R}^{d\times d}.
  \label{eq:Xi-def}
\end{equation}

\section{Discrete Evolution Law and Nonlinearity}

We consider the discrete-time dynamical system
\begin{equation}
  X_{t+1} = \Xi X_t + m T(X_t),
  \label{eq:evolution}
\end{equation}
where:

\begin{itemize}[leftmargin=*]
  \item \(\Xi\) is as in \eqref{eq:Xi-def},
  \item \(m\in\mathbb{R}\) is a scalar coupling, and
  \item \(T: \mathbb{R}^d\to\mathbb{R}^d\) is a nonlinearity applied entrywise.
\end{itemize}

In the main instantiation, \(T\) is given by
\[
  T(x) = (\tanh(x_1),\dots,\tanh(x_d))^\top,
\]
with global Lipschitz constant \(L_T = 1\).
A small fixed coupling \(m\) is used, e.g.\ \(m = 0.01\).

\subsection{Lipschitz Bound for the Nonlinearity}

\begin{lemma}[Lipschitz Constant of \(\tanh\)]
Let \(\tanh: \mathbb{R}\to\mathbb{R}\) denote the hyperbolic tangent.
Then
\[
  |\tanh(a) - \tanh(b)| \le |a-b|, \quad \forall a,b\in\mathbb{R}.
\]
Consequently, the vector map \(T: \mathbb{R}^d\to\mathbb{R}^d\),
\(T(x)_k = \tanh(x_k)\), is Lipschitz with \(L_T = 1\):
\[
  \|T(x) - T(y)\|_2 \le \|x-y\|_2, \quad \forall x,y\in\mathbb{R}^d.
\]
\end{lemma}

\begin{proof}
The derivative \(\tanh'(z) = \operatorname{sech}^2(z)\) satisfies
\(0 < \tanh'(z) \le 1\) for all \(z\in\mathbb{R}\). The mean value theorem implies
\(|\tanh(a)-\tanh(b)| \le \sup_z |\tanh'(z)| \cdot |a-b| \le |a-b|\).
The vector case follows by summing squared components and taking square roots.
\end{proof}

\section{Contraction Mapping Analysis}

Define the one-step map \(G: \mathbb{R}^d\to\mathbb{R}^d\) by
\begin{equation}
  G(x) = \Xi x + m T(x).
\end{equation}

\begin{lemma}[Lipschitz Bound for \(G\)]
\label{lem:G-Lipschitz}
Let \(\Xi\in\mathbb{R}^{d\times d}\), \(m\in\mathbb{R}\), and
\(T\) Lipschitz with constant \(L_T\). Then for all \(x,y\in\mathbb{R}^d\),
\[
  \|G(x) - G(y)\|_2
  \le (F + c)
       \|x-y\|_2,
\]
where \(F = \|\Xi\|_{\rm op}\) and \(c = |m|L_T\).
\end{lemma}

\begin{proof}
We have
\[
  G(x) - G(y) = \Xi(x-y) + m(T(x) - T(y)).
\]
By the triangle inequality and operator norm submultiplicativity,
\begin{align*}
  \|G(x)-G(y)\|_2
  &\le \|\Xi(x-y)\|_2 + |m|\,\|T(x)-T(y)\|_2 \\
  &\le \|\Xi\|_{\rm op}\,\|x-y\|_2 + |m|L_T\,\|x-y\|_2 \\
  &= (F + c)\,\|x-y\|_2.
\end{align*}
\end{proof}

\begin{theorem}[Sufficient Condition for Contraction]
\label{thm:contraction}
Under the assumptions of Lemma~\ref{lem:G-Lipschitz}, if
\[
  F + c < 1,
\]
then \(G\) is a contraction on \((\mathbb{R}^d,\|\cdot\|_2)\) with
contraction constant \(\kappa = F + c < 1\). There exists a unique fixed point
\(x^*\in\mathbb{R}^d\) such that \(G(x^*) = x^*\), and for every
initial condition \(x_0\in\mathbb{R}^d\) the iterates \(x_{t+1}=G(x_t)\)
converge exponentially to \(x^*\):
\[
  \|x_t - x^*\|_2 \le \kappa^t \|x_0 - x^*\|_2.
\]
\end{theorem}

\begin{proof}
Lemma~\ref{lem:G-Lipschitz} gives \(\|G(x)-G(y)\|_2 \le (F+c)\|x-y\|_2\).
If \(F + c < 1\), then \(G\) is a strict contraction on the complete
metric space \((\mathbb{R}^d,\|\cdot\|_2)\). The Banach Fixed Point
Theorem yields the existence and uniqueness of a fixed point and exponential
convergence with rate \(\kappa = F + c\).
\end{proof}

In the concrete setting with \(T(x) = \tanh(x)\) and constant \(m\), we take
\(L_T = 1\) and \(c = |m|\). The condition \(\|\Xi\|_{\rm op} + |m| < 1\)
is thus the key theorem-level criterion.

\section{Operator Norm Bounds for Aggregated Prime Families}

\subsection{Basic Triangular Bound}

For \(\Xi\) defined in \eqref{eq:Xi-def} we have
\begin{align}
  \|\Xi\|_{\rm op}
  &= \biggl\|\operatorname{Re}\Bigl(\sum_{p\in\mathcal{P}} w_p U_p\Bigr)\biggr\|_{\rm op} \\
  &\le \biggl\|\sum_{p\in\mathcal{P}} w_p U_p\biggr\|_{\rm op}
     \quad\text{(since \(\|\operatorname{Re}(A)\|_{\rm op} \le \|A\|_{\rm op}\))} \\
  &\le \sum_{p\in\mathcal{P}} |w_p| \|U_p\|_{\rm op}
     \quad\text{(triangle inequality)}.
\end{align}
If each \(U_p\) is unitary (so \(\|U_p\|_{\rm op}=1\)), this reduces to
\begin{equation}
  F = \|\Xi\|_{\rm op} \le \sum_{p\in\mathcal{P}} |w_p|.
  \label{eq:F-upper}
\end{equation}

\subsection{Empirical Saturation in Frozen Families}

In the frozen protocol used as a negative control, with
\(d=24\), \(\mathcal{P} = \{2,3\}\), and
\(w_p = \Omega(p)/p^{3/2}\) (so \(\sum_{p} |w_p| \approx 1.2531\)),
we consider three families:

\begin{enumerate}[leftmargin=*]
  \item Diagonal: \(U_p = D_p\).
  \item Fourier-conjugated: \(U_p = F^\dagger D_p F\).
  \item Permutation-conjugated: \(U_p = P_p^\dagger D_p P_p\).
\end{enumerate}

Numerical singular value decompositions for each family yield
\[
  F = \|\Xi\|_{\rm op} = \sum_{p\in\mathcal{P}} |w_p|
  \quad\text{(to within machine precision)}.
\]
Thus the basic upper bound \eqref{eq:F-upper} is saturated in these
families: no non-trivial operator norm cancellation is observed.
This empirically falsifies the hypothesis that ``nice'' basis mixing
(diagonal \(\to\) Fourier \(\to\) permutation) would automatically reduce
\(F\) below the naive weight sum.

\section{Optimised Families and Explicit Contraction}

\subsection{Optimisation Objective}

Motivated by the negative control, we treat the operator norm as an
explicit objective. In a Fourier-conjugated family, we parametrise
phases and allow signed/complex weights:
\begin{align}
  U_p(\theta_p) &= F^\dagger D_p(\theta_p) F, \\
  D_p(\theta_p) &= \operatorname{diag}(e^{i\theta_{p,0}},\dots,e^{i\theta_{p,d-1}}),
\end{align}
with \(\theta_p \in \mathbb{R}^d\) and constraint \(\|U_p(\theta_p)\|_{\rm op}=1\).
Weights \(w_p\in\mathbb{C}\) are allowed to be signed/complex but
regularised toward \(\Omega(p)/p^\alpha\).

We define
\begin{equation}
  \Xi^* = \operatorname{Re}\Bigl(\sum_{p\in\mathcal{P}} w_p U_p(\theta_p)\Bigr),
\end{equation}
with spectral norm \(F^* = \|\Xi^*\|_{\rm op}\), and solve
\begin{equation}
  \min_{\{\theta_p\},\{w_p\}} \; F^* 
  = \min \bigl\|\operatorname{Re}\bigl(\sum_p w_p U_p(\theta_p)\bigr)\bigr\|_{\rm op},
\end{equation}
subject to the unitary and regularisation constraints.

\subsection{Contraction Once \(F^* + |m| < 1\)}

When the optimisation produces a family with
\begin{equation}
  F^* + |m| < 1,
\end{equation}
Theorem~\ref{thm:contraction} applies directly with \(\Xi\) replaced by
\(\Xi^*\), yielding:

\begin{corollary}[Explicit Contraction in Optimised Families]
Let \(\Xi^*\) and \(F^*\) be as above, with \(T(x)=\tanh(x)\)
and constant \(m\in\mathbb{R}\). If \(F^* + |m| < 1\) then
\[
  G^*(x) = \Xi^* x + m T(x)
\]
 is a contraction on \(\mathbb{R}^d\) with contraction constant
\(\kappa^* = F^* + |m| < 1\), and admits a unique fixed point
\(x^*\) to which all trajectories converge exponentially.
\end{corollary}

\section{Illustrative Code Snippet}

The following Python-like pseudocode illustrates the computation of \(F\)
for a given family in the frozen protocol:

\begin{lstlisting}[language=Python,basicstyle=\ttfamily\small]
import numpy as np

def prime_weights(primes, alpha=1.5):
    # Baseline w_p = 1 / p**alpha (Omega(p) = 1 for primes)
    return {p: 1.0 / (p ** alpha) for p in primes}

def diagonal_D(p, d):
    k = np.arange(d)
    phases = np.exp(2j * np.pi * k / p)
    return np.diag(phases)

def fourier_matrix(d):
    k = np.arange(d)
    n = k[:, None]
    F = np.exp(-2j * np.pi * n * k / d) / np.sqrt(d)
    return F

def build_Xi(primes, d, mode="diagonal"):
    w = prime_weights(primes)
    if mode == "fourier":
        F = fourier_matrix(d)
    Xi = np.zeros((d, d), dtype=complex)
    for p in primes:
        Dp = diagonal_D(p, d)
        if mode == "diagonal":
            Up = Dp
        elif mode == "fourier":
            Up = F.conj().T @ Dp @ F
        else:
            raise ValueError("unsupported mode")
        Xi += w[p] * Up
    Xi_real = Xi.real  # Re(sum_p w_p U_p)
    # Operator norm via largest singular value
    svals = np.linalg.svd(Xi_real, compute_uv=False)
    F_norm = float(svals[0])
    return F_norm

primes = [2, 3]
F_diag = build_Xi(primes, d=24, mode="diagonal")
F_four = build_Xi(primes, d=24, mode="fourier")
print("F_diag =", F_diag)
print("F_four =", F_four)
\end{lstlisting}

This snippet encodes the core computation behind the negative control
experiments; optimisation extends it by parameterising phases and weights
and minimising \texttt{F\_norm} under the desired constraints.

\section{Prior Art Scope}

This report is intended as a defensive publication.
It discloses as prior art:

\begin{itemize}[leftmargin=*]
  \item The combination of prime-indexed unitary families \(\{U_p\}\),
        real-part aggregation \(\operatorname{Re}(\sum_p w_p U_p)\), and
        a small-gain style contraction criterion \(\|\Xi\|_{\rm op} + |m|L_T < 1\)
        in a discrete-time nonlinear system.
  \item The explicit use of negative controls (diagonal/Fourier/permutation
        families) to falsify naive interference hypotheses by showing
        \(\|\Xi\|_{\rm op} = \sum |w_p|\) to numerical precision.
  \item The optimisation programme over phases and signed/complex \(w_p\)
        constrained by a governance envelope (ADRs) that distinguishes
        theorem-level claims from empirical behaviour.
\end{itemize}

It does not attempt to claim standard contraction theorems, generic spectral
norm optimisation, or general prime number theory; these are well documented
in the existing literature and referenced in the accompanying bibliography.

\appendix
\section*{Mathematical Appendix}
\addcontentsline{toc}{section}{Mathematical Appendix}

\section{Preliminaries and Notation}

\subsection{Spaces and Norms}

Let \(d \in \mathbb{N}\) be the state dimension. We work on
\(\mathbb{R}^d\) (or \(\mathbb{C}^d\) when convenient) with the Euclidean norm
\[
  \|x\|_2 = \Bigl(\sum_{k=1}^d |x_k|^2\Bigr)^{1/2}.
\]
For a matrix \(A \in \mathbb{C}^{d\times d}\) we use:

\begin{itemize}[leftmargin=*]
  \item The operator (spectral) norm
  \[
    \|A\|_{\rm op}
    = \sup_{\|x\|_2 = 1} \|Ax\|_2
    = \sigma_{\max}(A),
  \]
  where \(\sigma_{\max}(A)\) is the largest singular value.
  \item The entrywise real part \(\operatorname{Re}(A)\), defined by
  \((\operatorname{Re}(A))_{ij} = \operatorname{Re}(A_{ij})\).
\end{itemize}

\subsection{Prime-Indexed Operators}

Let \(\mathcal{P}\) be a finite set of primes. For each \(p \in \mathcal{P}\)
and dimension \(d\), we define:

\begin{itemize}[leftmargin=*]
  \item \textbf{Diagonal prime-phase operator}
  \[
    D_p \in \mathbb{C}^{d\times d}, \qquad (D_p)_{kk} = e^{2\pi i k/p},
    \quad k = 0,\dots,d-1.
  \]
  \item \textbf{Fourier-conjugated operator} (for a fixed unitary DFT matrix
  \(F \in \mathbb{C}^{d\times d}\) with \(F^\dagger F = I_d\)):
  \[
    U_p^{\rm (F)} = F^\dagger D_p F.
  \]
  \item \textbf{Permutation-conjugated operator} (for a permutation matrix
  \(P_p \in \mathbb{R}^{d\times d}\)):
  \[
    U_p^{\rm (perm)} = P_p^\dagger D_p P_p.
  \]
\end{itemize}

In all cases above, the \(U_p\) are unitary, so
\(\|U_p\|_{\rm op} = 1\).

\subsection{Weights and the Universal Multiplicity Constant}

The Universal Multiplicity Constant \(\Lambda_m\) is not a scalar in this
appendix; rather, it serves as a conceptual label for a \emph{regime} of
prime-indexed weights and operators in which multiplicity (recurrence of
structure across primes) is stable under perturbation.

Concretely, we introduce prime-indexed weights
\(\{w_p\}_{p \in \mathcal{P}}\) of the form
\begin{equation}
  w_p = \frac{\Omega(p)}{p^\alpha},
  \label{eq:wp-def}
\end{equation}
where:
\begin{itemize}[leftmargin=*]
  \item \(\Omega(p)\) is a multiplicity function (e.g.\ the number of
  prime factors counted with multiplicity in extended models; for primes
  themselves we take \(\Omega(p) = 1\)),
  \item \(\alpha > 1\) is a decay exponent ensuring absolute summability
        when \(\mathcal{P}\) is extended.
\end{itemize}

Given \(\{U_p\}_{p\in\mathcal{P}}\) and \(\{w_p\}_{p\in\mathcal{P}}\), we define
\begin{equation}
  \Xi = \operatorname{Re}\Bigl(\sum_{p\in\mathcal{P}} w_p U_p\Bigr)
  \in \mathbb{R}^{d\times d}.
  \label{eq:Xi-appendix}
\end{equation}
The pair \((\{U_p\},\{w_p\})\) lies in the \emph{\(\Lambda_m\)-stable band}
if it satisfies the contraction inequality developed below.

\section{Discrete Evolution Law and Nonlinearity}

We consider the discrete-time dynamical system
\begin{equation}
  X_{t+1} = \Xi X_t + m\, T(X_t),
  \label{eq:evolution-appendix}
\end{equation}
where:

\begin{itemize}[leftmargin=*]
  \item \(\Xi\) is as in \eqref{eq:Xi-appendix};
  \item \(m \in \mathbb{R}\) is a scalar coupling (in applications \(m\)
  is small, e.g.\ \(m = 0.01\));
  \item \(T : \mathbb{R}^d \to \mathbb{R}^d\) is a nonlinearity applied
  entrywise.
\end{itemize}

In the main text, we fix
\[
  T(x) = (\tanh(x_1),\dots,\tanh(x_d))^\top.
\]

\subsection{Lipschitz Constant of the Nonlinearity}

\begin{lemma}[Lipschitz Constant of \(\tanh\)]
The function \(\tanh : \mathbb{R} \to \mathbb{R}\) is globally Lipschitz
with constant \(L_T = 1\), i.e.
\[
  |\tanh(a) - \tanh(b)| \le |a-b|, \quad \forall a,b \in \mathbb{R}.
\]
Consequently, the induced map
\(T : \mathbb{R}^d \to \mathbb{R}^d\),
\(T(x)_k = \tanh(x_k)\), satisfies
\[
  \|T(x) - T(y)\|_2 \le \|x-y\|_2, \quad \forall x,y\in\mathbb{R}^d.
\]
\end{lemma}

\begin{proof}
We have \(\tanh'(z) = \operatorname{sech}^2(z) \in (0,1]\) for all
\(z\in\mathbb{R}\), with \(\sup_z |\tanh'(z)| = 1\).
By the mean value theorem,
\(|\tanh(a)-\tanh(b)| \le \sup_z |\tanh'(z)|\,|a-b| \le |a-b|\).

For the vector map,
\begin{align*}
  \|T(x) - T(y)\|_2^2
  &= \sum_{k=1}^d |\tanh(x_k) - \tanh(y_k)|^2 \\
  &\le \sum_{k=1}^d |x_k - y_k|^2
   = \|x-y\|_2^2,
\end{align*}
and taking square roots yields the claim.
\end{proof}

\section{Contraction Mapping Structure}

Define the one-step map \(G : \mathbb{R}^d \to \mathbb{R}^d\) by
\begin{equation}
  G(x) = \Xi x + m\, T(x).
\end{equation}

\subsection{Lipschitz Bound for \(G\)}

\begin{lemma}[Lipschitz Constant of \(G\)]
\label{lem:G-Lipschitz}
Let \(\Xi \in \mathbb{R}^{d\times d}\), \(m\in\mathbb{R}\), and assume
\(T\) is Lipschitz with constant \(L_T\). Then for all \(x,y\in\mathbb{R}^d\),
\begin{equation}
  \|G(x) - G(y)\|_2
  \le (F + c)\,\|x-y\|_2,
\end{equation}
where \(F = \|\Xi\|_{\rm op}\) and \(c = |m|L_T\).
\end{lemma}

\begin{proof}
We compute
\[
  G(x) - G(y) = \Xi(x-y) + m\,(T(x) - T(y)).
\]
By the triangle inequality and submultiplicativity of \(\|\cdot\|_{\rm op}\),
\begin{align*}
  \|G(x) - G(y)\|_2
  &\le \|\Xi(x-y)\|_2 + |m|\,\|T(x) - T(y)\|_2 \\
  &\le \|\Xi\|_{\rm op}\,\|x-y\|_2 + |m| L_T\,\|x-y\|_2 \\
  &= (F + c)\,\|x-y\|_2.
\end{align*}
\end{proof}

\subsection{Sufficient Condition for Contraction}

\begin{theorem}[Sufficient Contraction Condition]
\label{thm:contraction}
Under the assumptions of Lemma~\ref{lem:G-Lipschitz}, if
\begin{equation}
  F + c < 1, \qquad F = \|\Xi\|_{\rm op},\quad c = |m|L_T,
\end{equation}
then \(G\) is a contraction on \((\mathbb{R}^d,\|\cdot\|_2)\) with
contraction constant \(\kappa = F + c < 1\). In particular:

\begin{itemize}[leftmargin=*]
  \item There exists a unique fixed point \(x^\ast\in\mathbb{R}^d\)
        such that \(G(x^\ast) = x^\ast\).
  \item For any initial condition \(x_0\in\mathbb{R}^d\), the iterates
  \(x_{t+1} = G(x_t)\) satisfy
  \[
    \|x_t - x^\ast\|_2 \le \kappa^t\,\|x_0 - x^\ast\|_2.
  \]
\end{itemize}
\end{theorem}

\begin{proof}
By Lemma~\ref{lem:G-Lipschitz},
\(\|G(x)-G(y)\|_2 \le (F + c)\,\|x-y\|_2\) for all \(x,y\).
If \(F + c < 1\), then \(G\) is a strict contraction on the complete
metric space \((\mathbb{R}^d,\|\cdot\|_2)\). The Banach Fixed Point
Theorem yields uniqueness of the fixed point and exponential convergence
with rate \(\kappa = F + c\).
\end{proof}

In the concrete setup with \(T(x) = \tanh(x)\) and constant \(m\),
we have \(L_T = 1\) and hence \(c = |m|\). The Universal Multiplicity
Constant \(\Lambda_m\) governs the selection of \(\{U_p\},\{w_p\}\)
so that \(\|\Xi\|_{\rm op} + |m| < 1\) can be achieved in a structured
prime-indexed manner.

\section{Operator Norm Bounds and the Role of \(\Lambda_m\)}

We now connect the operator norm \(F = \|\Xi\|_{\rm op}\) to the
prime-indexed structure.

\subsection{Triangle Inequality Bound}

For \(\Xi\) defined in \eqref{eq:Xi-appendix} we have:
\begin{align}
  \|\Xi\|_{\rm op}
  &= \biggl\|\operatorname{Re}\Bigl(\sum_{p\in\mathcal{P}} w_p U_p\Bigr)\biggr\|_{\rm op} \\
  &\le \biggl\|\sum_{p\in\mathcal{P}} w_p U_p\biggr\|_{\rm op}
     \quad\text{(since \(\|\operatorname{Re}(A)\|_{\rm op} \le \|A\|_{\rm op}\))} \\
  &\le \sum_{p\in\mathcal{P}} |w_p|\,\|U_p\|_{\rm op}
     \quad\text{(triangle inequality and homogeneity)}.
\end{align}
If each \(U_p\) is unitary (or satisfies \(\|U_p\|_{\rm op}\le 1\)),
this simplifies to
\begin{equation}
  F = \|\Xi\|_{\rm op}
  \le \sum_{p\in\mathcal{P}} |w_p|.
  \label{eq:F-basic-bound}
\end{equation}

\subsection{Empirical Saturation in Frozen Families}

In the frozen protocol used as a negative control, with
\(d = 24\), \(\mathcal{P} = \{2,3\}\), and
\(w_p = \Omega(p)/p^{3/2}\) (so \(\sum_{p} |w_p| \approx 1.2531\)),
we consider three families:

\begin{enumerate}[leftmargin=*]
  \item Diagonal: \(U_p = D_p\).
  \item Fourier-conjugated: \(U_p = F^\dagger D_p F\).
  \item Permutation-conjugated: \(U_p = P_p^\dagger D_p P_p\).
\end{enumerate}

Numerical SVD computations yield:
\[
  F = \|\Xi\|_{\rm op} = \sum_{p\in\mathcal{P}} |w_p|
  \quad\text{(to within machine precision)}.
\]
Thus the upper bound \eqref{eq:F-basic-bound} is saturated: no
non-trivial operator-norm cancellation occurs in these families.
This falsifies a naive ``nice interference'' hypothesis in which
non-commuting prime-phase mixing was expected to reduce \(F\)
automatically.

From the perspective of \(\Lambda_m\), these families lie outside the
\(\Lambda_m\)-stable band: even though they are aesthetically
prime-structured, they fail the contraction criterion
\(\|\Xi\|_{\rm op} + |m| < 1\).

\section{Optimised Prime-Structured Families and \(\Lambda_m\)-Stability}

\subsection{Optimisation Objective}

To exhibit \(\Lambda_m\)-stability explicitly, we formulate
an optimisation problem over phase parameters and weights in a
Fourier-conjugated family:
\begin{align}
  U_p(\theta_p) &= F^\dagger D_p(\theta_p) F, \\
  D_p(\theta_p) &= \operatorname{diag}(e^{i\theta_{p,0}},\dots,e^{i\theta_{p,d-1}}),
\end{align}
with \(\theta_p \in \mathbb{R}^d\) and \(\|U_p(\theta_p)\|_{\rm op}=1\)
for all \(p\).

Weights \(w_p \in \mathbb{C}\) are allowed to be signed/complex but
regularised toward the baseline pattern \eqref{eq:wp-def},
so that the multiplicity semantics of \(\Lambda_m\) are preserved.

We define
\begin{equation}
  \Xi^\star = \operatorname{Re}\Bigl(\sum_{p\in\mathcal{P}} w_p U_p(\theta_p)\Bigr),
\end{equation}
with operator norm \(F^\star = \|\Xi^\star\|_{\rm op}\), and solve
\begin{equation}
  \min_{\{\theta_p\},\{w_p\}} \; F^\star
  = \min \bigl\|\operatorname{Re}\bigl(\sum_p w_p U_p(\theta_p)\bigr)\bigr\|_{\rm op},
\end{equation}
subject to the unitary and regularisation constraints.

A solution \((\{\theta_p\},\{w_p\})\) is said to lie in the
\emph{\(\Lambda_m\)-stable band} if
\begin{equation}
  F^\star + |m| < 1.
  \label{eq:lambda-m-stable}
\end{equation}

\subsection{Contraction in the \(\Lambda_m\)-Stable Band}

When \eqref{eq:lambda-m-stable} holds, Theorem~\ref{thm:contraction}
applies directly:

\begin{corollary}[Contraction in a \(\Lambda_m\)-Stable Family]
Let \(\Xi^\star\) and \(F^\star\) be as above, with \(T(x) = \tanh(x)\)
and constant \(m\). If
\[
  F^\star + |m| < 1,
\]
then the map
\[
  G^\star(x) = \Xi^\star x + m\,T(x)
\]
is a contraction on \((\mathbb{R}^d,\|\cdot\|_2)\) with contraction
constant \(\kappa^\star = F^\star + |m| < 1\), admits a unique fixed
point \(x^\star\), and all trajectories \(x_{t+1} = G^\star(x_t)\)
converge exponentially to \(x^\star\).
\end{corollary}

\subsection{Robustness Considerations}

In practice, \(\Lambda_m\)-stability is assessed numerically and then
interpreted structurally:

\begin{itemize}[leftmargin=*]
  \item Numerically, one verifies \(F^\star + |m| < 1\) for the
        optimised family in several dimensions and for multiple initial
        conditions, observing uniform convergence.
  \item Structurally, one then analyses the pattern of
        \(\{w_p\},\{\theta_p\}\) to identify recurrent relations across
        primes that characterise the \(\Lambda_m\)-stable band, so that
        stability is not an isolated artefact but a multiplicity pattern.
\end{itemize}

In a more formal setting, one could use interval analysis or perturbation
bounds to certify that the inequality \(F^\star + |m| < 1\) persists
under small deviations in \((\{\theta_p\},\{w_p\})\).

\section{Summary of Bounds and Logical Structure}

For convenience, we summarise the key inequalities and logical steps
used in the article:

\begin{itemize}[leftmargin=*]
  \item \textbf{Nonlinear Lipschitz bound:}
    \[
      \|T(x) - T(y)\|_2 \le L_T \|x-y\|_2, \quad \text{with } L_T = 1
      \text{ for } T(x)=\tanh(x).
    \]
  \item \textbf{Aggregate Lipschitz bound:}
    \[
      \|G(x) - G(y)\|_2 \le (\|\Xi\|_{\rm op} + |m|L_T)\,\|x-y\|_2.
    \]
  \item \textbf{Sufficient contraction condition:}
    if \(\|\Xi\|_{\rm op} + |m|L_T < 1\) then \(G\) is a contraction,
    with unique fixed point and exponential convergence.
  \item \textbf{Triangle inequality bound for \(\Xi\):}
    \[
      \|\Xi\|_{\rm op}
      \le \sum_{p\in\mathcal{P}} |w_p|\,\|U_p\|_{\rm op},
      \quad \text{so } \|\Xi\|_{\rm op} \le \sum_{p} |w_p|
      \text{ when } \|U_p\|_{\rm op}\le 1.
    \]
  \item \textbf{Empirical saturation in frozen families:}
    for diagonal/Fourier/permutation families with fixed \(\mathcal{P}\)
    and \(d\), numerical SVD shows
    \(\|\Xi\|_{\rm op} = \sum_{p} |w_p|\), so the basic bound is tight
    and no cancellation occurs.
  \item \textbf{Optimised \(\Lambda_m\)-stable families:}
    by optimising phases and weights in a prime-indexed family, one
    constructs examples for which \(F^\star + |m| < 1\), and the
    contraction theorem applies with explicit constant
    \(\kappa^\star = F^\star + |m|\).
\end{itemize}

From the perspective of the Universal Multiplicity Constant \(\Lambda_m\),
the negative results demonstrate that merely labeling operators by primes
is not sufficient for stability: only certain prime-indexed patterns of
phases and weights produce the operator norm bounds required for
contraction. Those patterns comprise the \(\Lambda_m\)-stable band and
serve as the backbone for the multiplicity-theoretic interpretation in
the main text.

\nocite{*}
\bibliographystyle{unsrt}  
\bibliography{references}  


\end{document}
